% mage-auto-heading: accessibility-and-the-engineering-experiment
% mage-auto-heading: what-the-user-sees
% mage-auto-heading: adding-a-description-without-losing-the-document
% mage-auto-heading: one-system-several-engineering-questions
% mage-auto-heading: following-one-job
% mage-auto-heading: how-the-book-uses-docable
% Authoritative LaTeX manuscript; Markdown is a generated compatibility projection.

% <!-- part-title: The New Engineering Problem -->

% <!-- chapter-title: MAGE by Example: Summary of the DocAble Case -->

\chapter{MAGE by Example: Summary of the DocAble Case}

% <!-- point: teaser-grounds-the-abstractions | You will meet DocAble in every section; here is just enough to hold it. | terms: running-example -->

\href{\%7B\%7Bproduct_url\%7D\%7D}{DocAble} provides the book\textquotesingle s deepest worked example: a production accessibility system built largely through coding agents. \magevalue{chapter:software-factories} presents the full case study; before examining its models and Alignment mechanisms here, a few facts about the system are in order.

\section{Accessibility and the Engineering Experiment}\label{heading:accessibility-and-the-engineering-experiment}

% <!-- section-terms: document-accessibility, supervised-autonomy -->

% <!-- point: explicit-problem-inaccessible-documents-at-scale | The explicit problem: inaccessible university documents at scale. | terms: document-accessibility, remediation-cost -->

In 2024, the U.S. Department of Justice gave public universities a deadline to make their digital content accessible, measured against the Web Content Accessibility Guidelines (\magevalue{doj_technical_standard}). A typical slide deck can fail the standard in dozens of ways: images without descriptions, missing titles, or a reading order that a screen reader cannot follow. Purdue's manual remediation effort made the scale of the work concrete. Over approximately six months, \magevalue{seip_uta_workers} trained undergraduate teaching assistants billed \magevalue{seip_uta_hours} hours; the resulting corpus contained \magevalue{seip_uta_pairs} identifiable before-and-after document pairs from \magevalue{seip_uta_courses} courses. That is about \magevalue{seip_uta_hours_per_pair} billed hours per identifiable pair.\footnote{The recovered pairs need not account for every document the students worked on. This corpus-level ratio is not a measured remediation time for each document.}\citep{davis2026docable}

I personally maintain two courses and had no desire to spend a full semester remediating them manually. I could delegate some of that work to other people, but delegation would not make me any less responsible for the result. I was, however, perfectly willing to spend six months attempting to automate the work to my own standard.

% <!-- point: two-birds-one-stone-motive | One project answered both problems: build the production system, and use the engineering work to discover what sustained agentic development demanded. | terms: document-accessibility, running-example -->

In early 2026, coding agents had also become capable enough to make a different experiment possible. I wanted to know what software engineering would look like if I stopped writing most of the implementation myself, so I used accessibility remediation as a real production problem on which to find out. The system became DocAble; the resulting method became MAGE. MAGE makes the judgment of one engineer capable of governing the implementation capacity of a software team. \magevalue{chapter:software-factories} returns to the build and examines how that capability emerged.

\section{What the User Sees}\label{heading:what-the-user-sees}

% <!-- section-terms: governance-target-product, document-accessibility -->

% <!-- point: upload-in-remediated-file-plus-evidence-out | From the user's view DocAble is simple: an inaccessible document goes in; a remediated document and an evidence record come back. | terms: governance-target-product -->

% <!-- index-def: fidelity-validator -->

% <!-- index-def: provenance-layer -->

From the user\textquotesingle s perspective, DocAble is simple: upload an inaccessible document; receive a remediated document and an evidence record. It works across common document formats: Word, PowerPoint, and Excel; their LibreOffice equivalents; PDF; and text-based formats including LaTeX and Typst. Figures are described, reading order is repaired, and structural accessibility is restored where possible.

Why not simply give the entire document to a capable multimodal model and ask it to make the document accessible? I tried that too. For documents within their operating range, capable models can produce impressive results. But the approach is expensive because the model must reconstruct properties already available deterministically in the document (the document\textquotesingle s title may be right there on the page, but represented as large bold text rather than Word\textquotesingle s Title style); large documents can exceed its practical operating range (the model may simply give up before finishing the deck); and support varies across formats (a model that handles PowerPoint may not accept PDF).

More importantly, "the model said it was accessible" is not an adequate compliance case. To trust that result, we would still have to make explicit the obligations against which it should be judged. The governing obligation is not merely to produce a plausible-looking artifact, but to provide students with a substantially equivalent experience. If a result is challenged, the university needs credible evidence of what the system checked (every image was examined for a usable description), what it changed (this image received this alternative text), what it could establish (the repaired reading order matches the intended sequence), and what limitations remain (this element could not be remediated automatically). DocAble therefore treats accessibility remediation as an engineering problem rather than a single act of inference.

% <!-- point: the-figure-stops-at-the-boundary | The glimpse stops at the boundary; before looking inside, we need only know what engineering result the system is supposed to produce. | terms: governance-target-product -->

Before looking inside the system, \autoref{docable-glimpse} shows only the engineering result it is supposed to produce.

\begin{figure}[htbp]
\centering
\includegraphics[width=\linewidth]{book/latex/assets/docable-glimpse-d34d027fb5.pdf}
\caption{\emph{DocAble from the user\textquotesingle s view.} An inaccessible document goes in; a remediated document and evidence of consequential changes come back. Internal structure is omitted because it is not yet relevant.}
\label{docable-glimpse}
\end{figure}

% <!-- point: the-interface-hides-a-large-production-system | The simple interface hides a large production system. | terms: governance-target-product -->

That simple interface hides a large production system. At the final snapshot of the 31-week record, October 14, 2026, DocAble contained \magevalue{seip_prod_loc} lines of production code and \magevalue{seip_apparatus_loc} lines of engineering apparatus, about 3.8 times as much. The apparatus includes tests, checks, tools, infrastructure, models, and indexed documentation. The examples that follow isolate individual engineering questions from that system. Together they show different places where Modeling and Alignment enter the work.

% <!-- point: four-kinds-of-obligation-state-the-case | What DocAble had to accomplish falls into four kinds of obligation. | terms: governance-target-product -->

\autoref{docable-obligation-kinds} groups what that system had to accomplish into four kinds of obligation.

{\footnotesize
\begin{longtable}[]{@{}
  >{\raggedright\arraybackslash}p{(\linewidth - 2\tabcolsep) * \real{0.5000}}
  >{\raggedright\arraybackslash}p{(\linewidth - 2\tabcolsep) * \real{0.5000}}@{}}
\caption[Four kinds of obligation]{What DocAble had to accomplish, grouped into four kinds of obligation. \magevalue{chapter:software-factories} states the complete requirements this grouping summarizes.}\label{docable-obligation-kinds}\tabularnewline
\toprule\noalign{}
\begin{minipage}[b]{\linewidth}\raggedright
Kind of obligation
\end{minipage} & \begin{minipage}[b]{\linewidth}\raggedright
DocAble must\ldots{}
\end{minipage} \\
\midrule\noalign{}
\endfirsthead
\toprule\noalign{}
\begin{minipage}[b]{\linewidth}\raggedright
Kind of obligation
\end{minipage} & \begin{minipage}[b]{\linewidth}\raggedright
DocAble must\ldots{}
\end{minipage} \\
\midrule\noalign{}
\endhead
\bottomrule\noalign{}
\endlastfoot
Genuine accessibility & Repair problems that require understanding document content, judged against the governing standard (WCAG 2.1 AA) rather than proxy checker scores. \\
Preservation & Retain the document\textquotesingle s content, appearance, and behavior except where remediation requires a change. \\
Accountability & Record what motivated each change and what evidence supports the result, keeping consequential changes inspectable and, where appropriate, reversible. \\
Practical operation at scale & Handle the formats people actually use, at practical sizes, costs, and latency, as a deployed service for thousands of faculty. \\
\end{longtable}
}

% <!-- point: two-tensions-make-it-an-engineering-problem | Two tensions among the obligations make DocAble an engineering problem rather than a file-conversion task: accessibility can conflict with preservation, and semantic quality pulls against economical automation. | terms: document-accessibility, governance-target-product -->

Two tensions among these obligations make DocAble an engineering problem rather than a file-conversion task. Accessibility can conflict with preservation: repairing a reading order or describing a figure means changing the document, and the change must not damage the author\textquotesingle s content, appearance, or editability. And semantic quality pulls against economical automation: some repairs need the contextual interpretation foundation models supply, but model output is uncertain and can be expensive, so the system must combine that inference with deterministic processing and validation.

% <!-- point: examples-establish-the-problem-full-requirements-follow | The examples establish the nature of the engineering problem; the factories chapter presents the complete requirements and examines how they shaped the system. | terms: governance-target-product -->

These examples establish the nature of the engineering problem. \magevalue{chapter:software-factories} presents the complete requirements and examines how they shaped the resulting system.

\section{Adding a Description Without Losing the Document}\label{heading:adding-a-description-without-losing-the-document}

% <!-- point: a-useful-description-must-preserve-the-authors-document | Describing a lecture diagram requires interpretation in context, while the lecturer's original content must survive the repair. | terms: running-example, document-accessibility -->

Consider a lecture slide containing a diagram and several sentences explaining it. A student using a screen reader needs a description that conveys what the diagram contributes to the lesson. ``A diagram'' would be a description, but not a particularly useful one. Producing something better requires interpreting the figure in context. Meanwhile, the lecturer\textquotesingle s original content must survive. Making a document accessible is not permission to rewrite the lecture.

% <!-- point: structured-document-operations-make-the-repair-tractable | DocAble's structured document models and named operations expose the objects and changes relevant to a repair without repeatedly reconstructing the file format's machinery. | terms: running-example, modeling-principle -->

DocAble uses structured document models to represent the objects its repairs operate on. Remediation code works through named document operations rather than manipulating the underlying file format wherever convenient. This gives the engineering question a manageable form: which object needs a description, what information should be attached to it, and which existing content must remain? The representation exposes the parts relevant to the repair without requiring every repair to rediscover the document format\textquotesingle s machinery.

% <!-- point: preservation-requires-inclusion-rather-than-equality | Comparable content extracted from the original and revised documents lets DocAble require original content to remain while allowing added accessibility information. | terms: running-example, modeling-principle -->

Preservation requires another purposeful representation. DocAble extracts comparable content from the original and revised documents. The original extracted content must remain present in the result. Equality would be the wrong requirement: the purpose of the repair is to add accessibility information. The relationship must allow those additions while detecting the disappearance of content the author supplied.

% <!-- point: a-production-validator-enforces-the-preservation-condition | DocAble's production validator fails a job when the extracted-content preservation condition is violated; structured operations restrict repairs, and the validator checks their result. | terms: running-example, alignment-principle -->

Alignment gives that requirement a consequence. A production validator compares the extracted content before and after remediation and fails the job when the preservation condition is violated. A repair does not become acceptable merely because it runs successfully or reports that it has finished. The resulting document must pass the check. The structured operations restrict how repairs are made; the validator checks a property of what they produced.

% <!-- point: preservation-checks-leave-semantic-judgment-open | The foundation model can choose a description, but content inclusion establishes only survival of the extracted content, not meaningful alternative text or unchanged layout. | terms: running-example -->

There is still useful judgment to delegate. The foundation model can interpret the diagram and choose how to describe it without an engineer writing every description in advance. The preservation check does not establish that the description is meaningful, nor does content inclusion alone establish that the layout is unchanged. Those are separate properties requiring suitable evidence. What the check does establish is narrower and useful: the content represented by its extractor has survived.

% <!-- point: the-preservation-obligation-survives-implementation-choices | Coding agents may choose how to implement repairs within the document-model interface while the explicit preservation requirement and its enforcement continue to apply. | terms: running-example, modeling-principle, alignment-principle -->

This arrangement also guides the coding agents that build and change DocAble. They can choose how to implement a new repair within the document-model interface, while the preservation requirement continues to apply. The engineer has made a consequential obligation explicit and given it an enforcement mechanism. Much of the implementation remains open. That combination---preserving what matters while leaving room to solve the problem---is the opportunity MAGE develops.

\section{One System, Several Engineering Questions}\label{heading:one-system-several-engineering-questions}

% <!-- section-terms: running-example -->

% <!-- point: building-well-is-several-problems | Building DocAble well is not one engineering problem; different questions require different representations. | terms: running-example -->

Building DocAble well is not one engineering problem. Different questions require different representations.

% <!-- point: five-questions-five-representations | Structure, behavior, decisions, measurement, and provenance are different questions, each wanting its own representation of the same system. | terms: running-example -->

Ask what the system contains and how its parts relate, and you want a structural view. Ask what happens to a document after upload, and you want behavior: the states a job moves through and the transitions that are legal. Ask which component may call another, and you want a model of permitted relationships. Ask whether the system is fast and cheap enough, and you want measurements and their allowable envelopes. Ask whether the evidence returned with a document faithfully describes what happened, and you want provenance.

% <!-- point: no-single-model-of-docable | The system has not changed between these questions; the representation has. No single model is best for all of them — each is a purposeful reduction that preserves the facts one question needs and suppresses detail it does not. | terms: running-example, model-selection -->

The system has not changed between these questions. The representation has. There is no single model of DocAble that is best for all of them. Each model is a purposeful reduction that preserves the facts needed for an engineering question and suppresses detail that question does not need.

\section{Following One Job}\label{heading:following-one-job}

% <!-- section-terms: running-example -->

% <!-- point: enough-common-territory-to-locate-the-models | Before Chapters 2 and 3 examine particular properties of DocAble, one ordinary job through the system shows what those properties concern. | terms: running-example -->

Before Chapters \magevalue{chapter:modeling|num} and \magevalue{chapter:alignment|num} examine particular properties of DocAble, \autoref{one-docable-job} follows one ordinary job through the system.

\begin{figure}[htbp]
\centering
\includegraphics[width=\linewidth]{book/latex/assets/one-docable-job-bbd94ab765.pdf}
\caption{\emph{Processing one document in DocAble.} A user submits a document, creating a job that tracks its processing. The work may be divided into chunks and processed by workers using specialized services for generation, rendering, optical character recognition (OCR), and document-format processing. The resulting changes are reassembled and validated, and the remediated document is returned with evidence.}
\label{one-docable-job}
\end{figure}

% <!-- index-def: remediation-job -->

% <!-- index-def: worker -->

% <!-- index-def: chunk -->

% <!-- index-def: remediation-service -->

% <!-- index-def: work-queue -->

% <!-- index-def: evidence-record -->

% <!-- point: the-territory-vocabulary | A job tracks one document's trip; a chunk is a portion of that work processed independently; workers claim chunks, calling specialized services; queues coordinate handoffs; the outputs are the remediated document and its evidence. | terms: running-example -->

A job tracks one document\textquotesingle s trip through the system. A chunk is a portion of that work that can be processed independently. Workers claim and process chunks, calling specialized services for capabilities such as image description, rendering, conversion, or OCR. Queues and other infrastructure coordinate those handoffs. The final outputs are the remediated document and evidence describing consequential changes and remaining findings.

% <!-- point: one-worker-several-models | The same worker appears differently in later models — a deployable component, an actor in a lifecycle, an owner of a resource, a subject of an access policy, or a source of measurements; not competing descriptions, but different properties of the same system. | terms: running-example, model-selection -->

The same worker will appear differently in later models. It may be a deployable component, an actor in a lifecycle, an owner of a resource, a subject of an access policy, or a source of measurements. These are not competing descriptions. Each preserves different properties of the same system for a different engineering question: structure, behavior, ownership, decisions, measurements, or provenance.

\section{How the Book Uses DocAble}\label{heading:how-the-book-uses-docable}

% <!-- section-terms: running-example -->

% <!-- point: docable-shown-at-method-resolution | DocAble is the book's deepest worked example because its engineering can be shown at the resolution the method requires; Chapter 2 asks which representations make properties tractable, Chapter 3 how intent comes to bind the work. | terms: running-example, modeling-principle, alignment-principle -->

DocAble is the book\textquotesingle s deepest worked example because the system can be examined at the resolution the method requires. Chapters \magevalue{chapter:modeling|num} and \magevalue{chapter:alignment|num} use the working system to ask two different kinds of questions. \magevalue{chapter:modeling} asks which representations make important system properties tractable. \magevalue{chapter:alignment} asks how engineering intent comes to bind the work.

% <!-- point: motion-then-the-production-system | Chapter 4 examines how those mechanisms are used during development, following engineering problems through successive models, decisions, implementation, and governance conversion; Chapter 5 then examines the production system itself — what failed, what those failures revealed, and how recurring judgments became durable engineering structure. | terms: running-example -->

\magevalue{chapter:one-problem-many-models} examines how those mechanisms are used during development. It follows engineering problems through successive models, decisions, implementation, and governance conversion, showing how Modeling and Alignment change as the system changes. \magevalue{chapter:software-factories} then examines the production system itself: what failed, what those failures revealed, and how recurring judgments became models, controls, and other durable engineering structure.

% <!-- point: docable-is-depth-not-the-whole-evidence | DocAble is the primary case study; Chapter 5 examines practices at other engineering organizations, identifying where related structures recur, where they take different forms, and where the evidence supports a more limited interpretation. | terms: running-example -->

DocAble is the primary case study. \magevalue{chapter:software-factories} examines practices at other engineering organizations, identifying where related structures recur, where they take different forms, and where the available evidence supports more limited conclusions.

% <!-- point: bridge-to-the-substrate | We can now examine the machinery that makes autonomous implementation possible and the engineering mechanisms through which MAGE acts on it. | terms: running-example -->

We can now examine the machinery that makes autonomous implementation possible and the engineering mechanisms through which MAGE acts on it.
